\section{与前几期的连接：Agent、AI Native 与模型公司}

EP128 与前几集构成一组很清晰的 Agent 产品谱系。EP139 讲 Agent 的技术史和产品化漏斗；EP130 从张月光角度讲 AI Native、One Way Door 和 Agent 产品心智；EP129 从智谱角度讲模型公司、Agent 作为模型落地路径；EP128 则展示一个 Agent 产品如何真的进入 ARR 和收购事件。

\subsection{四期的不同侧面}

本节把四期放在一张表里，方便读者看到同一个 Agent 问题的不同剖面。

\begin{longtable}{p{0.18\textwidth}p{0.36\textwidth}p{0.37\textwidth}}
\toprule
期数 & 侧重点 & 和 Manus 的关系 \\
\midrule
EP139 & Agent 技术史、评测、产品化漏斗 & 提供 Manus 所在赛道的技术背景。 \\
EP130 & AI Native 产品、One Way Door、AI 朋友 & 解释 Manus 为什么要守住用户心智和单向门。 \\
EP129 & 模型公司、Agent 落地、开源闭源 & 解释 Manus 为什么依赖模型生态和基础设施。 \\
EP128 & Manus 增长、制造业隐喻、复杂性风险 & 给出一个真实 Agent 公司案例。 \\
\bottomrule
\end{longtable}

\begin{importantbox}{系列中的 Manus}
Manus 是 Agent 从概念走向商业化的样本。它不是技术史的终点，也不是 AI Native 的唯一答案，但它把“用户真的用、真的付费、真的产生工作价值”这件事推到了前台。
\end{importantbox}

\subsection{本章小结}

EP128 的意义在于把前几集的抽象概念落成一个具体公司：Agent 的技术、产品、组织、商业化和复杂性风险，在 Manus 身上同时出现。

